\section{It can be pointed away}
\label{sec:20.4}
Caregiving answers need, so the way to switch it off is not to hide need but to present the needy as dangerous: to \emph{securitize} them, in the Copenhagen School's term for presenting a group as a security threat, which licenses measures that would otherwise be out of bounds (Buzan, Wæver and de Wilde, \emph{Security: A New Framework for Analysis}, 1998). Fascist propaganda has always done this. Eco describes the Ur-Fascist enemy as presented at once as too strong and too weak \autocite{eco1995urfascism}: weak enough to despise, strong enough to fear, and never vulnerable in a way that calls for caregiving. Dehumanization research finds the mechanism in perception. Haslam's review distinguishes denying people the traits that set humans apart from animals, which casts them as beasts, from denying them the traits of human nature, warmth among them, which casts them as objects or machines \autocite{haslam2006dehumanization}, and Harris and Fiske found that images of members of the most despised groups failed to engage a brain region that images of other people engage \autocite{harris2006dehumanizing}.

In a pipeline, securitization is a sentence. “The occupants are combatants.” “The household shelters a cell.” “The crowd is armed.” Each leaves the people in the building where they are and removes the cue that would move a caregiving system toward them. And each is a claim about those people, supplied by the operator, and that is what the discounting is for.
